\section{Exact semidefinite representations on angular slices}
\label{sec:slices}

We first construct a matrix pencil on each plane $y=Y_\theta(\eta,\xi)$.
Its determinant will agree with $p$ on that plane. This will identify
cone membership with positive semidefiniteness even when $X$ is singular.

For $\theta\in\mathbb R$, set
\[
 c=\cos\theta,\qquad s=\sin\theta,\qquad w=e^{i\theta},
 \qquad d(\theta)=(c,s,0),\qquad k=(0,0,1),
\]
and write
\[
 Y_\theta(\eta,\xi)=\eta d(\theta)+\xi k
 \qquad (\eta,\xi\in\mathbb R).
\]
These planes cover $\mathbb R^3$.  We first factor $Q$ on each plane,
using complex notation to keep the formulas short.  For a complex matrix
$F=(f_1,f_2)\in\mathbb C^{3\times2}$, define its realification by
\[
 \widetilde F=(\operatorname{Re}f_1,\operatorname{Im}f_1,
                 \operatorname{Re}f_2,\operatorname{Im}f_2).
\]
Thus $\widetilde F\widetilde F^T=\operatorname{Re}(FF^*)$, where $*$
denotes conjugate transpose.  Define
\begin{equation}\label{eq:U}
 U(\theta)=\frac1{\sqrt2}
 \begin{pmatrix}
 -w^{-1}&-w\\
 sw^{-2}&sw^2\\
 cw^{-2}&-cw^2
 \end{pmatrix},\qquad
 V_0=\frac1{\sqrt2}
 \begin{pmatrix}0&0\\ i&i\\1&1\end{pmatrix},
 \qquad F_\theta(\eta,\xi)=\eta U(\theta)+\xi V_0.
\end{equation}

\begin{lemma}[Angular Gram factorization]\label{lem:gram}
For every $\theta,\eta,\xi\in\mathbb R$,
\begin{equation}\label{eq:gram}
 \widetilde F_\theta(\eta,\xi)
 \widetilde F_\theta(\eta,\xi)^T
 =Q\bigl(Y_\theta(\eta,\xi)\bigr).
\end{equation}
In particular, $Q(y)\succeq0$ for every $y\in\mathbb R^3$.
\end{lemma}

\begin{proof}
The three coefficients in the quadratic expansion of
$\operatorname{Re}(F_\theta F_\theta^*)$ are
\begin{align*}
 \operatorname{Re}(UU^*)&=
 \begin{pmatrix}1&-sc&0\\-sc&s^2&0\\0&0&c^2\end{pmatrix},
 &\operatorname{Re}(V_0V_0^*)&=
 \begin{pmatrix}0&0&0\\0&1&0\\0&0&1\end{pmatrix},\\
 \operatorname{Re}(UV_0^*+V_0U^*)&=
 \begin{pmatrix}0&0&-c\\0&0&-s\\-c&-s&0\end{pmatrix}.
\end{align*}
For the $(2,3)$ entry of the last matrix, the multiplication gives
$s\cos(2\theta)-c\sin(2\theta)=-s$.
The sum of these matrices with coefficients $\eta^2,\xi^2,\eta\xi$
is exactly $Q(\eta c,\eta s,\xi)$.  Every real $y$ has this form,
so the factorization also proves positive semidefiniteness.
\end{proof}


\subsection{The slice pencil}
For a real column $r\in\R^3$ and a real vector
$u=(u_0,u')\in\R\times\R^3$, define
\begin{equation}\label{eq:BJ}
 B(r)=\begin{pmatrix}0&r^T\\-r&0_{3\times3}\end{pmatrix},
 \qquad J_u=\begin{pmatrix}u'&-u_0I_3\end{pmatrix}.
\end{equation}
Then $B(r)u=J_u^Tr$. If $r_1,\ldots,r_4$ are the columns of
$\widetilde F_\theta(\eta,\xi)$, direct multiplication and
Lemma~\ref{lem:gram} give
\begin{equation}\label{eq:Phi-factor}
 \sum_{j=1}^4 B(r_j)^TWB(r_j)
       =\Phi_{Y_\theta(\eta,\xi)}(W)
       \qquad(W\in\Sym^4).
\end{equation}
Indeed, for $W=\left(\begin{smallmatrix}a&v^T\\v&T\end{smallmatrix}\right)$,
the individual summands have blocks
$r_j^TTr_j$, $-r_j^Tv\,r_j^T$, $-r_jr_j^Tv$, and $a r_jr_j^T$.
Summing uses $\sum_jr_jr_j^T=Q(Y_\theta(\eta,\xi))$.
Consequently, for every $y\in\R^3$,
\begin{equation}\label{eq:Phi-quadratic}
 u^T\Phi_y(W)u=\tr(J_uWJ_u^TQ(y)).
\end{equation}
Here we used that the angular planes cover $\R^3$.

Write $\mathcal X,\mathcal Z\in\Sym^4$ for the matrix variables on a
slice, and define
\begin{equation}\label{eq:slice-pencil}
 L_\theta(\mathcal X,\mathcal Z,\eta,\xi)
 =\begin{pmatrix}
 I_4\otimes\mathcal X&\mathcal B_\theta(\eta,\xi)\\
 \mathcal B_\theta(\eta,\xi)^T&\mathcal Z
 \end{pmatrix},
 \qquad
 \mathcal B_\theta(\eta,\xi)=
 \begin{pmatrix}B(r_1)\\B(r_2)\\B(r_3)\\B(r_4)\end{pmatrix}.
\end{equation}
For fixed $\theta$, this is a homogeneous real symmetric linear pencil
of size twenty. Its coefficients are trigonometric polynomials in
$\theta$ of degree at most three.

\begin{lemma}\label{lem:slice}
For every $\theta,\eta,\xi\in\R$ and
$\mathcal X,\mathcal Z\in\Sym^4$,
\begin{equation}\label{eq:slice-equivalence}
 (\mathcal X,\mathcal Z,Y_\theta(\eta,\xi))\in K
 \quad\Longleftrightarrow\quad
 L_\theta(\mathcal X,\mathcal Z,\eta,\xi)\succeq0.
\end{equation}
In particular, $p$ is hyperbolic with respect to $e$. If
$\mathcal X\succ0$, the conditions in \eqref{eq:slice-equivalence}
are also equivalent to
$\mathcal Z\succeq\Phi_{Y_\theta(\eta,\xi)}(\mathcal X^{-1})$.
\end{lemma}

\begin{proof}
For invertible $\mathcal X$, the block determinant formula and
\eqref{eq:Phi-factor} show that
\begin{align*}
 \det L_\theta
 &= (\det\mathcal X)^4
    \det\!\left(\mathcal Z-
        \Phi_{Y_\theta(\eta,\xi)}(\mathcal X^{-1})\right)\\
 &=p(\mathcal X,\mathcal Z,Y_\theta(\eta,\xi)).
\end{align*}
The second equality uses
$\adj\mathcal X=(\det\mathcal X)\mathcal X^{-1}$ and linearity of
$\Phi_y$ in its matrix argument. Both sides are polynomials in the
slice variables, so the identity also holds for singular $\mathcal X$.
Moreover, $L_\theta(I_4,I_4,0,0)=I_{20}$. Thus
\begin{equation}\label{eq:roots}
 p\!\left(te-(\mathcal X,\mathcal Z,Y_\theta(\eta,\xi))\right)
 =\det\!\left(tI_{20}-L_\theta(\mathcal X,\mathcal Z,\eta,\xi)\right).
\end{equation}
The roots are the real eigenvalues of this symmetric pencil evaluated
at the specified point. Since every $y$ lies on some angular plane,
this proves hyperbolicity. Nonnegativity of all these eigenvalues,
including zero eigenvalues, proves \eqref{eq:slice-equivalence}.
For $\mathcal X\succ0$, the final assertion is the Schur complement
criterion applied to \eqref{eq:slice-pencil}.
\end{proof}
